\section{Approche ADMM associée aux ondelettes biorthogonales 4.4}

\subsection{Problème inverse en tomodensitométrie}
\begin{frame}{Modélisation du problème}
        \begin{equation}
            \mathbf{g} = A \mathbf{f} + \mathbf{n}
        \end{equation}
        où \begin{itemize}
            \item \( \mathbf{f} \) : image (coefficients d'atténuation)
            \item \( \mathbf{g} \) : projections (sinogramme)
            \item \( A \) : transformée de Radon discrète
            \item \( \mathbf{n} \) : bruit
        \end{itemize}
\end{frame}

\subsection{Mal-positude}
\begin{frame}{Mal-positude}
    \transfade
    \begin{block}<1->{Difficulté de la reconstruction}
        \begin{itemize}
            \item Sous-déterminé
            \item Sensible au bruit
            \item Non unicité des solutions
        \end{itemize}
    \end{block}

    \begin{block}<2->{Conséquences}
        \begin{itemize}
            \item artefacts
            \item instabilité
            \item mauvaise qualité image
        \end{itemize}
    \end{block}
\end{frame}

\subsection{Ondelettes orthogonales}
\begin{frame}{Ondelettes orthogonales}

    \textbf{Principe fondamental}
    \begin{itemize}
        \item La base d'analyse et la base de synthèse coïncident
        \item L'opérateur de transformation $W$ est une matrice orthogonale
    \end{itemize}

    \textbf{Propriétés mathématiques}
    \begin{itemize}
        \item L'adjoint coïncide avec l'inverse :
        \[
        W^T = W^{-1}
        \]
        \item Orthogonalité des colonnes :
        \[
        W^T W = I
        \]
        \item Conservation de l'énergie (propriété de Parseval) :
        \[
        \|Wx\|_2 = \|x\|_2
        \]
        \item Représentation non redondante : $m = n$
    \end{itemize}
    % \textbf{Avantages et limitations}
    % \begin{itemize}
    %     \item[+] Conservation exacte de l'énergie
    %     \item[+] Reconstruction parfaite et simple ($W^{-1} = W^T$)
    %     \item[+] Transformation bijective et sans redondance
    %     \item[-] Filtres généralement non symétriques
    %     \item[-] Absence de phase linéaire (distorsion possible des contours)
    % \end{itemize}
    \textbf{Exemples :} Haar, Daubechies, Symlets

\end{frame}


\subsection{Ondelettes biorthogonales}
\begin{frame}{Ondelettes biorthogonales}
    \textbf{Principe fondamental}
    \begin{itemize}
        \item Deux bases distinctes :
        \begin{itemize}
            \item Base d'analyse : $W$ (décomposition)
            \item Base de synthèse : $\tilde{W}$ (reconstruction)
        \end{itemize}
        \item Filtres d'analyse : $h$, $g$ (passe-bas, passe-haut)
        \item Filtres de synthèse : $\tilde{h}$, $\tilde{g}$ (passe-bas, passe-haut)
    \end{itemize}

    \vspace{0.2cm}

    \textbf{Propriétés mathématiques}
    \begin{itemize}
        \item Reconstruction parfaite :
        \[
        \tilde{W} W = I
        \]
        \item L'adjoint de l'analyse n'est pas l'opérateur de synthèse :
        \[
        W^T \neq \tilde{W}
        \]
        \item Condition de biorthogonalité :
        \[
        \langle \psi_{j,k}, \tilde{\psi}_{j',k'} \rangle = \delta_{j,j'} \delta_{k,k'}
        \]
    \end{itemize}

    % \vspace{0.2cm}

    % \textbf{Avantages pratiques}
    % \begin{itemize}
    %     \item[+] Filtres symétriques possibles → phase linéaire
    %     \item[+] Meilleure préservation des contours et structures géométriques
    %     \item[+] Flexibilité accrue dans la conception des filtres
    %     \item[+] Qualité de reconstruction supérieure pour l'imagerie
    % \end{itemize}

    % \vspace{0.2cm}

    % \textbf{À retenir pour l'optimisation}
    % \begin{itemize}
    %     \item Reconstruction effective : $\tilde{W}$ (synthèse)
    %     \item Gradient / optimisation : $W^T$ (adjoint mathématique)
    %     \item \textcolor{red}{Attention :} Certaines implémentations utilisent $\tilde{W}$ à la place de $W^T$
    % \end{itemize}

    % \vspace{0.2cm}
    \textbf{Exemples :} Bior4.4, Bior2.2, Bior6.8

\end{frame}

\subsection{Reconstruction tomographique à faible dose par parcimonie: Compressed Sensing}
\begin{frame}{Modélisation CT et exploitation de la parcimonie}
    \transfade
    \only<1>{\begin{block}{Modèle d'acquisition CT}
        \begin{equation}
            \mathbf{g} = A \mathbf{f} \in \mathbb{R}^m, \quad A \in \mathbb{R}^{m \times n} \text{ (matrice système)}
        \end{equation}
        \begin{itemize}
            \item $\mathbf{f} \in \mathbb{R}^n$ : image à reconstruire
            \item $\mathbf{g} \in \mathbb{R}^m$ : projections (sinogramme)
            \item $A$ : modélise la géométrie CT
        \end{itemize}
    \end{block}}

    \only<2>{\begin{block}{Avec représentation parcimonieuse}
        \begin{equation}
            \mathbf{g} = A \mathbf{f} = A W \mathbf{x} = \Phi' \mathbf{x}, \quad \Phi' = A W
        \end{equation}
        \begin{itemize}
            \item $W$ : matrice d'ondelettes
            \item $\mathbf{x}$ : coefficients parcimonieux dans la base $W$
            \item $\Phi'$ : matrice combinée système CT + ondelettes
        \end{itemize}
        \vspace{0.5em}
        \centering
        \textit{On travaille maintenant directement avec les coefficients parcimonieux $\mathbf{x}$ pour reconstruire $\mathbf{f}$ à partir de peu de mesures.}
    \end{block}}

\end{frame}

\subsection{Approche ADMM combinée avec ondelettes (Bior4.4)}
\begin{frame}{ADMM combinée à bior4.4}
    \transfade
    \only<1>{\begin{block}{Principe ADMM}
        Décomposition du problème
        On introduit :
        \[
        z = W\mathbf{f}
        \]
        
        \[
        \min_{\mathbf{f},z}
        \frac{1}{2} ||A\mathbf{f} - \mathbf{g}||_2^2 + \lambda ||z||_1
        \]
    \end{block}}
    \only<2>{\begin{block}{Lagrangien augmenté}
        \[
        L_\rho(\mathbf{f},z,u) =
        \frac{1}{2} ||A\mathbf{f} - \mathbf{g}||_2^2
        + \lambda |z|_1
        + \frac{\rho}{2} ||W\mathbf{f} - z + u||_2^2
          \]
    \end{block}
    où \begin{itemize}
            \item $\mathbf{x}$ : variable primaire, vecteur des coefficients d'image à reconstruire
            \item $z$ : variable auxiliaire pour séparer le terme $\ell_1$ et faciliter l'optimisation
            \item $u$ : variable duale mise à l'échelle (dual scaled), associée à la contrainte $Wx = z$
            \item $W$ : matrice de transformation (par exemple, base d'ondelettes ou opérateur de gradient)
            \item $A$ : opérateur de mesure (matrice système ou projection CT)
            \item $\mathbf{y}$ : vecteur de mesures (sinogramme)
            \item $\lambda$ : paramètre de régularisation, contrôle l'importance de la parcimonie
            \item $\rho > 0$ : paramètre de pénalité, contrôle l'écart à la contrainte $Wx = z$ et influence la vitesse de convergence
        \end{itemize}
    }
    \only<3>{
    \begin{block}{Algorithme ADMM}
        \begin{itemize}
            \item \textbf{Mise à jour de $\mathbf{x}$ (sous-problème quadratique)}  
            Résolution du système :
            \[
            x^{k+1} = \arg\min_{x} \frac{1}{2}\|\mathcal{A}x - y\|_2^2 + \frac{\rho}{2}\|Wx - z^k + u^k\|_2^2
            \]
            ou équivalent linéaire :
            \[
            (A^T A + \rho W^T W)x = A^T y + \rho W^T (z^k - u^k)
            \]

            \item \textbf{Mise à jour de $z$ (seuillage doux)}  
            Imposition de la parcimonie via l’opérateur proximal $\ell_1$ :
            \[
            z^{k+1} = \mathcal{S}_{\lambda / \rho}(Wx^{k+1} + u^k), \quad
            \mathcal{S}_{\tau}(t) = \operatorname{sign}(t)\max(0, |t|-\tau)
            \]

            \item \textbf{Mise à jour de la variable duale $u$}  
            Correction de la violation de contrainte $Wx = z$ :
            \[
            u^{k+1} = u^k + (Wx^{k+1} - z^{k+1})
            \]
        \end{itemize}
    \end{block}}
\end{frame}
